On verse sur $\qty{16.8}{\g}$ de poudre de fer, un volume $V=\qty{250}{\mL}$
d'une solution d'acide chlorhydrique de concentration $c=\qty{2.0}{\mol\per\L}$.

Il se forme du dihydrogène et des ions fer(II) en solution.
\begin{enumerate}
    \item Écrire l'équation de la réaction en considérant l'acide chlorhydrique
          comme une solution d'ions \ch{H^+_{(aq)}}.
    \item Calculer les quantités initiales des réactifs.
    \item Quel est le réactif limitant~?
\end{enumerate}

